\section{Discussion}

The results presented in this work provide empirical evidence that the outer dynamics of galaxy discs are modulated by their large-scale environment. In particular, we find, using the SPARC sample, a statistically significant correlation between the environmental overdensity proxy $\delta_{\rm mass}$ and the outer rotation-curve observable $F_3$, with $\rho \approx 0.52$, together with a systematic improvement in model selection ($\Delta\mathrm{AIC} \approx -7.5$) when including the environmental term. Out-of-sample (OOS) validation further shows a robust reduction in predictive error, with median $\Delta\mathrm{RMSE}_{\rm out} \approx -7.27\ \mathrm{km\ s^{-1}}$ and consistent improvement across all validation splits.

These results support the inclusion of an effective environmental term in the phenomenological description of outer-disc kinematics. Importantly, this term is introduced in a model-agnostic manner, without assuming a specific theoretical framework (e.g.\ dark matter halo parametrizations or modified gravity prescriptions). Instead, the approach is entirely data-driven: the environmental dependence emerges empirically as an additional predictor that improves both in-sample model selection and out-of-sample predictive performance.

A key result of this work is that the environmental contribution remains significant after controlling for internal galaxy properties, including baryonic mass and size. This indicates that the observed correlation is not trivially reducible to internal structural parameters, but reflects an additional degree of freedom associated with the surrounding large-scale structure. In this sense, the SCM framework can be interpreted as a minimal phenomenological extension that captures this dependence without introducing strong theoretical assumptions.

The robustness of the environmental signal is supported by several independent checks. First, the improvement in predictive performance persists across randomized cross-validation splits, with 100\% of validation realizations yielding $\Delta\mathrm{RMSE}_{\rm out} < 0$. Second, the statistical significance of the effect is confirmed by a paired Wilcoxon test ($p \approx 1.19 \times 10^{-10}$). Third, the signal remains stable under bootstrap resampling and when using alternative environmental proxies (e.g.\ nearest-neighbour surface density), indicating that it is not driven by a particular subset of galaxies or by a specific definition of environment.

We emphasise that the environmental term introduced here is purely phenomenological and data-driven, and does not imply a specific underlying physical mechanism. Any physical interpretation of this term remains open and subject to further investigation. In particular, we do not claim a modification of gravity, nor do we propose a complete cosmological model. The present results should instead be interpreted as empirical evidence that the outer regions of galaxy discs are sensitive to environmental conditions, in a way that is not fully captured by models based solely on internal baryonic structure.

These findings have several implications. First, they suggest that outer-disc kinematics may provide a new observational probe of the galaxy--environment connection, complementary to traditional measures based on morphology or star formation. Second, they motivate the inclusion of environmental terms in future modelling efforts, including both semi-empirical approaches and numerical simulations. In particular, testing whether similar correlations arise in hydrodynamical simulations or large-scale structure models will be an important next step.

Finally, the framework presented here makes testable predictions. If the environmental modulation is real, one expects systematic differences in outer rotation-curve behaviour between galaxies residing in high-density environments (e.g.\ groups and clusters) and those in low-density regions (voids), even at fixed baryonic mass. Extending the present analysis to larger and more diverse samples, as well as to independent datasets, will be essential to further assess the generality of the effect.

In summary, the results presented here provide a statistically robust, reproducible indication that environmental effects contribute to shaping the outer dynamics of galaxy discs. The SCM framework offers a simple and testable way to quantify this contribution, while remaining agnostic about its physical origin.
